% Options for packages loaded elsewhere
% Options for packages loaded elsewhere
\PassOptionsToPackage{unicode}{hyperref}
\PassOptionsToPackage{hyphens}{url}
\PassOptionsToPackage{dvipsnames,svgnames,x11names}{xcolor}
%
\documentclass[sigconf]{webmedia}

\usepackage{lmodern}
\usepackage[brazil]{babel}
\usepackage{hyperref}

%% BibTeX logo
\AtBeginDocument{%
  \providecommand\BibTeX{{%
    \normalfont B\kern-0.5em{\scshape i\kern-0.25em b}\kern-0.8em\TeX}}}

\settopmatter{printacmref=false, printfolios=false}

%% Configuração do evento -- ajuste conforme necessário
\setevent{webmedia}
\proceedingsDetails[WebMedia'2026]{Proceedings of the Brazilian Symposium on Multimedia and the Web}{2026}{Lavras/MG, Brazil}
\ISSN{2966-2753}
\usepackage{xcolor}
\usepackage{amsmath,amssymb}
\setcounter{secnumdepth}{-\maxdimen} % remove section numbering
\usepackage{iftex}
\ifPDFTeX
  \usepackage[T1]{fontenc}
  \usepackage[utf8]{inputenc}
  \usepackage{textcomp} % provide euro and other symbols
\else % if luatex or xetex
  \usepackage{unicode-math} % this also loads fontspec
  \defaultfontfeatures{Scale=MatchLowercase}
  \defaultfontfeatures[\rmfamily]{Ligatures=TeX,Scale=1}
\fi
\usepackage{lmodern}
\ifPDFTeX\else
  % xetex/luatex font selection
\fi
% Use upquote if available, for straight quotes in verbatim environments
\IfFileExists{upquote.sty}{\usepackage{upquote}}{}
\IfFileExists{microtype.sty}{% use microtype if available
  \usepackage[]{microtype}
  \UseMicrotypeSet[protrusion]{basicmath} % disable protrusion for tt fonts
}{}
\makeatletter
\@ifundefined{KOMAClassName}{% if non-KOMA class
  \IfFileExists{parskip.sty}{%
    \usepackage{parskip}
  }{% else
    \setlength{\parindent}{0pt}
    \setlength{\parskip}{6pt plus 2pt minus 1pt}}
}{% if KOMA class
  \KOMAoptions{parskip=half}}
\makeatother
% Make \paragraph and \subparagraph free-standing
\makeatletter
\ifx\paragraph\undefined\else
  \let\oldparagraph\paragraph
  \renewcommand{\paragraph}{
    \@ifstar
      \xxxParagraphStar
      \xxxParagraphNoStar
  }
  \newcommand{\xxxParagraphStar}[1]{\oldparagraph*{#1}\mbox{}}
  \newcommand{\xxxParagraphNoStar}[1]{\oldparagraph{#1}\mbox{}}
\fi
\ifx\subparagraph\undefined\else
  \let\oldsubparagraph\subparagraph
  \renewcommand{\subparagraph}{
    \@ifstar
      \xxxSubParagraphStar
      \xxxSubParagraphNoStar
  }
  \newcommand{\xxxSubParagraphStar}[1]{\oldsubparagraph*{#1}\mbox{}}
  \newcommand{\xxxSubParagraphNoStar}[1]{\oldsubparagraph{#1}\mbox{}}
\fi
\makeatother


\usepackage{longtable,booktabs,array}
\usepackage{calc} % for calculating minipage widths
% Correct order of tables after \paragraph or \subparagraph
\usepackage{etoolbox}
\makeatletter
\patchcmd\longtable{\par}{\if@noskipsec\mbox{}\fi\par}{}{}
\makeatother
% Allow footnotes in longtable head/foot
\IfFileExists{footnotehyper.sty}{\usepackage{footnotehyper}}{\usepackage{footnote}}
\makesavenoteenv{longtable}
\usepackage{graphicx}
\makeatletter
\newsavebox\pandoc@box
\newcommand*\pandocbounded[1]{% scales image to fit in text height/width
  \sbox\pandoc@box{#1}%
  \Gscale@div\@tempa{\textheight}{\dimexpr\ht\pandoc@box+\dp\pandoc@box\relax}%
  \Gscale@div\@tempb{\linewidth}{\wd\pandoc@box}%
  \ifdim\@tempb\p@<\@tempa\p@\let\@tempa\@tempb\fi% select the smaller of both
  \ifdim\@tempa\p@<\p@\scalebox{\@tempa}{\usebox\pandoc@box}%
  \else\usebox{\pandoc@box}%
  \fi%
}
% Set default figure placement to htbp
\def\fps@figure{htbp}
\makeatother


% definitions for citeproc citations
\NewDocumentCommand\citeproctext{}{}
\NewDocumentCommand\citeproc{mm}{%
  \begingroup\def\citeproctext{#2}\cite{#1}\endgroup}
\makeatletter
 % allow citations to break across lines
 \let\@cite@ofmt\@firstofone
 % avoid brackets around text for \cite:
 \def\@biblabel#1{}
 \def\@cite#1#2{{#1\if@tempswa , #2\fi}}
\makeatother
\newlength{\cslhangindent}
\setlength{\cslhangindent}{1.5em}
\newlength{\csllabelwidth}
\setlength{\csllabelwidth}{3em}
\newenvironment{CSLReferences}[2] % #1 hanging-indent, #2 entry-spacing
 {\begin{list}{}{%
  \setlength{\itemindent}{0pt}
  \setlength{\leftmargin}{0pt}
  \setlength{\parsep}{0pt}
  % turn on hanging indent if param 1 is 1
  \ifodd #1
   \setlength{\leftmargin}{\cslhangindent}
   \setlength{\itemindent}{-1\cslhangindent}
  \fi
  % set entry spacing
  \setlength{\itemsep}{#2\baselineskip}}}
 {\end{list}}
\usepackage{calc}
\newcommand{\CSLBlock}[1]{\hfill\break\parbox[t]{\linewidth}{\strut\ignorespaces#1\strut}}
\newcommand{\CSLLeftMargin}[1]{\parbox[t]{\csllabelwidth}{\strut#1\strut}}
\newcommand{\CSLRightInline}[1]{\parbox[t]{\linewidth - \csllabelwidth}{\strut#1\strut}}
\newcommand{\CSLIndent}[1]{\hspace{\cslhangindent}#1}



\setlength{\emergencystretch}{3em} % prevent overfull lines

\providecommand{\tightlist}{%
  \setlength{\itemsep}{0pt}\setlength{\parskip}{0pt}}



 


\makeatletter
\@ifpackageloaded{caption}{}{\usepackage{caption}}
\AtBeginDocument{%
\ifdefined\contentsname
  \renewcommand*\contentsname{Table of contents}
\else
  \newcommand\contentsname{Table of contents}
\fi
\ifdefined\listfigurename
  \renewcommand*\listfigurename{List of Figures}
\else
  \newcommand\listfigurename{List of Figures}
\fi
\ifdefined\listtablename
  \renewcommand*\listtablename{List of Tables}
\else
  \newcommand\listtablename{List of Tables}
\fi
\ifdefined\figurename
  \renewcommand*\figurename{Figure}
\else
  \newcommand\figurename{Figure}
\fi
\ifdefined\tablename
  \renewcommand*\tablename{Table}
\else
  \newcommand\tablename{Table}
\fi
}
\@ifpackageloaded{float}{}{\usepackage{float}}
\floatstyle{ruled}
\@ifundefined{c@chapter}{\newfloat{codelisting}{h}{lop}}{\newfloat{codelisting}{h}{lop}[chapter]}
\floatname{codelisting}{Listing}
\newcommand*\listoflistings{\listof{codelisting}{List of Listings}}
\makeatother
\makeatletter
\makeatother
\makeatletter
\@ifpackageloaded{caption}{}{\usepackage{caption}}
\@ifpackageloaded{subcaption}{}{\usepackage{subcaption}}
\makeatother
\makeatletter
\@ifpackageloaded{tcolorbox}{}{\usepackage[skins,breakable]{tcolorbox}}
\makeatother
\makeatletter
\@ifundefined{shadecolor}{\definecolor{shadecolor}{rgb}{.97, .97, .97}}{}
\makeatother
\makeatletter
\makeatother
\makeatletter
\ifdefined\Shaded\renewenvironment{Shaded}{\begin{tcolorbox}[enhanced, sharp corners, frame hidden, borderline west={3pt}{0pt}{shadecolor}, interior hidden, breakable, boxrule=0pt]}{\end{tcolorbox}}\fi
\makeatother
\usepackage{bookmark}
\IfFileExists{xurl.sty}{\usepackage{xurl}}{} % add URL line breaks if available
\urlstyle{same}
\hypersetup{
  pdftitle={Título do Trabalho},
  pdfauthor={Nome do Autor 1; Nome do Autor 2},
  pdfkeywords={palavra-chave1, palavra-chave2, palavra-chave3},
  colorlinks=true,
  linkcolor={blue},
  filecolor={Maroon},
  citecolor={Blue},
  urlcolor={Blue},
  pdfcreator={LaTeX via pandoc}}


\author{Nome do Autor 1}
\email{autor1@instituicao.br}
\affiliation{%
  \institution{Nome da Instituição}
  \city{Cidade}
  \state{}
  \country{Brasil}
}
\author{Nome do Autor 2}
\email{autor2@instituicao.br}
\affiliation{%
  \institution{Nome da Instituição}
  \city{Cidade}
  \state{}
  \country{Brasil}
}

\renewcommand{\shortauthors}{Autor1 et al.}

\title{Título do Trabalho}
\subtitle{Subtítulo (opcional)}
\begin{document}
\maketitle

\begin{abstract}
Insira aqui o resumo do trabalho.
\end{abstract}

\keywords{palavra-chave1, palavra-chave2, palavra-chave3}


\section{Cálculo de Determinação do Grafo a ser
Utilizado}\label{cuxe1lculo-de-determinauxe7uxe3o-do-grafo-a-ser-utilizado}

Descreva aqui o processo utilizado para determinar qual grafo será
analisado no trabalho. Apresente os critérios de seleção, as
características dos dados originais e justifique a escolha com base nos
objetivos do estudo.

O código-fonte, scripts e conjunto de dados utilizados estão disponíveis
no repositório:

\begin{itemize}
\tightlist
\item
  \textbf{Repositório GitHub:}
  \url{https://github.com/seu-usuario/seu-repositorio}
\item
  Códigos/scripts utilizados (Python ou R; se Python, recomenda-se
  Jupyter Notebook: \url{https://jupyter.org/})
\item
  Conjunto de dados tratados
\item
  \texttt{README} com instruções de execução do código
\end{itemize}

\section{Descrição do Grafo
Original}\label{descriuxe7uxe3o-do-grafo-original}

Apresente as características do grafo original: número de vértices,
número de arestas, tipo (dirigido/não-dirigido,
ponderado/não-ponderado), densidade, componentes conexas, distribuição
de grau, entre outros atributos relevantes.

\subsection{Tratamento Realizado}\label{tratamento-realizado}

Descreva os pré-processamentos e tratamentos aplicados ao grafo antes da
análise. Exemplos incluem: remoção de nós isolados, normalização de
pesos, filtragem de arestas duplicadas, conversão de formato, tratamento
de valores ausentes, entre outros. Mencione as ferramentas e scripts
utilizados em cada etapa.

\section{Metodologia}\label{metodologia}

Descreva os métodos, algoritmos e técnicas aplicados na análise do
grafo. Inclua referências às métricas calculadas (e.g., centralidade de
grau, betweenness, closeness, coeficiente de agrupamento, pagerank) e
justifique as escolhas metodológicas em relação aos objetivos do
trabalho.

Exemplo de equação de centralidade de grau:

\[
C_D(v) = \frac{\deg(v)}{n - 1}
\]

onde \(\deg(v)\) é o grau do vértice \(v\) e \(n\) é o número total de
vértices no grafo.

\section{Resultados Obtidos}\label{resultados-obtidos}

Apresente os resultados da análise. Utilize tabelas e gráficos para
ilustrar os achados. Para cada visualização, inclua um parágrafo
explicativo.

\subsection{Tabelas e Gráficos}\label{tabelas-e-gruxe1ficos}

Insira aqui as tabelas e gráficos com os resultados obtidos. Exemplo de
tabela com métricas do grafo:

\begin{longtable}[]{@{}ll@{}}
\caption{Métricas gerais do grafo
analisado}\label{tbl-metricas}\tabularnewline
\toprule\noalign{}
\textbf{Métrica} & \textbf{Valor} \\
\midrule\noalign{}
\endfirsthead
\toprule\noalign{}
\textbf{Métrica} & \textbf{Valor} \\
\midrule\noalign{}
\endhead
\bottomrule\noalign{}
\endlastfoot
Número de vértices & 0000 \\
Número de arestas & 0000 \\
Densidade & 0.000 \\
Diâmetro & 00 \\
Grau médio & 0.00 \\
Coef. de agrupamento & 0.000 \\
\end{longtable}

\begin{figure}

\centering{

\begin{verbatim}
[Distribuição de Grau - Substitua pela figura real]
\end{verbatim}

}

\caption{\label{fig-dist-grau}Distribuição de grau do grafo.
{[}Substitua pela figura real.{]}}

\end{figure}%

\subsection{Comentários para Cada
Tabela/Gráfico}\label{comentuxe1rios-para-cada-tabelagruxe1fico}

\textbf{Table~\ref{tbl-metricas}:} Apresente aqui uma interpretação dos
valores encontrados. Discuta o que cada métrica revela sobre a estrutura
do grafo, comparando com os valores esperados para o tipo de rede em
estudo.

\textbf{Figure~\ref{fig-dist-grau}:} Descreva o padrão observado na
distribuição de grau. Verifique se a distribuição segue uma lei de
potência (scale-free), distribuição aleatória ou outro modelo.

\subsection{Respostas às Perguntas
Obrigatórias}\label{respostas-uxe0s-perguntas-obrigatuxf3rias}

Responda objetivamente cada uma das perguntas obrigatórias propostas na
atividade:

\begin{enumerate}
\def\labelenumi{\arabic{enumi}.}
\tightlist
\item
  \textbf{{[}Pergunta 1{]}:} Insira aqui a pergunta e sua resposta
  fundamentada nos resultados obtidos.
\item
  \textbf{{[}Pergunta 2{]}:} Insira aqui a pergunta e sua resposta
  fundamentada nos resultados obtidos.
\item
  \textbf{{[}Pergunta N{]}:} Insira aqui a pergunta e sua resposta
  fundamentada nos resultados obtidos.
\end{enumerate}

\section{Discussão Crítica}\label{discussuxe3o-cruxedtica}

Interprete os resultados obtidos de forma crítica. Relacione os achados
com a literatura existente sobre análise de redes e grafos (Autor 2026).
Discuta:

\begin{itemize}
\tightlist
\item
  Limitações da análise realizada;
\item
  Possíveis vieses introduzidos pelo processo de tratamento dos dados;
\item
  Comparação com resultados de trabalhos relacionados;
\item
  Implicações práticas ou teóricas dos achados.
\end{itemize}

\section{Conclusão}\label{conclusuxe3o}

Sintetize os principais achados do trabalho, retomando os objetivos
iniciais e verificando se foram alcançados. Destaque as contribuições do
estudo e aponte possíveis direcionamentos para trabalhos futuros, como a
análise de grafos dinâmicos, a aplicação de técnicas de aprendizado de
máquina sobre o grafo, ou a extensão da metodologia para outros
domínios.

\subsection*{Agradecimentos}\label{agradecimentos}
\addcontentsline{toc}{subsection}{Agradecimentos}

Insira aqui os agradecimentos a pessoas, instituições ou fontes de
financiamento que contribuíram para a realização do trabalho.

\section*{Apêndice}\label{apuxeandice}
\addcontentsline{toc}{section}{Apêndice}

\subsection{Repositório e Instruções de
Execução}\label{reposituxf3rio-e-instruuxe7uxf5es-de-execuuxe7uxe3o}

O código-fonte completo está disponível em:
\url{https://github.com/seu-usuario/seu-repositorio}

O repositório contém:

\begin{itemize}
\tightlist
\item
  \textbf{scripts/} --- códigos Python/R utilizados nas análises
  (recomenda-se Jupyter Notebook: \url{https://jupyter.org/})
\item
  \textbf{data/} --- conjunto de dados tratados utilizados
\item
  \textbf{README.md} --- instruções detalhadas de como executar o código
  e reproduzir os resultados
\end{itemize}

\phantomsection\label{refs}
\begin{CSLReferences}{1}{0}
\bibitem[\citeproctext]{ref-exemplo-ref}
Autor, Exemplo. 2026. {``Título Do Artigo de Exemplo.''} \emph{Revista
de Exemplo} 1: 1--10.

\end{CSLReferences}




\end{document}
